%-----------------------------------------------------------------------

\documentclass[a4paper]{article}

\usepackage[utf8]{inputenc}
\usepackage[brazil]{babel}
\usepackage[T1]{fontenc}
\usepackage{inconsolata}
\usepackage{mathpazo}

\usepackage{indentfirst, graphicx, calc, setspace, icomma}
\usepackage[hmargin={1.7cm, 1.7cm}, vmargin={1.7cm, 1.7cm}]{geometry}
\usepackage{amsmath, amsfonts, mathrsfs, amssymb, amsthm} % , amsxtra
\usepackage[hang]{caption}

%-----------------------------------------------------------------------
% Suporte ao R.

$if(highlighting-macros)$
$highlighting-macros$
$endif$

$if(graphics)$
\usepackage{graphicx}
\makeatletter
\def\maxwidth{
  \ifdim\Gin@nat@width>\linewidth\linewidth\else\Gin@nat@width\fi}
\def\maxheight{
  \ifdim\Gin@nat@height>\textheight\textheight\else\Gin@nat@height\fi}
\makeatother
\setkeys{Gin}{width=\maxwidth, height=\maxheight, keepaspectratio}
$endif$

%-----------------------------------------------------------------------
% Definições feitas no template default para Latex do Rmarkdown.

% Necessário para listas.
\providecommand{\tightlist}{%
  \setlength{\itemsep}{0pt}\setlength{\parskip}{0pt}}

%-----------------------------------------------------------------------
% Tikz.

\usepackage{tikz}
\usepackage{pgfplots}
\usepackage{pgfplotstable}

\usetikzlibrary{arrows, calc, matrix, positioning, trees,
                decorations.pathmorphing, backgrounds, fit}

%-----------------------------------------------------------------------
% Definições.

% Linhas horizontais.
\newcommand{\HRule}{\noindent\rule{\linewidth}{0.2mm}}

%-----------------------------------------------------------------------

\begin{document}
\pagestyle{empty}

\HRule\\
\noindent
\begin{minipage}[t]{0.75\textwidth}
  \baselineskip 12pt
  $if(chair)$$chair$$if(class)$ - $class$$endif$\\ $endif$
  $if(lecturer)$ \hfill $url$\\ $endif$
  $if(lecturer)$$lecturer$$if(email)$ ($email$)$endif$\\ $endif$
  $if(department)$$department$\\ $endif$
  $title$ $if(date)$\hfill $date$$endif$\\
  Acadêmico: \hfill grr: \hspace{2cm}
\end{minipage}
\hfill
\begin{minipage}[t]{0.25\textwidth}
  \begin{flushright}
    \vspace*{-0.35cm}
    \includegraphics[height=2cm]{ufpr_logo.png}
  \end{flushright}
\end{minipage}\\
\HRule

\hspace{0.5cm}

$body$

\end{document}
